% ============================================================
\section{Conclusion}
\label{sec:conclusion}

We have specified, analyzed, and characterized M\textsuperscript{2}, a tokenomic primitive that combines a fixed-supply redemption-backed token, an asymmetric-fee Uniswap-v4 hook, and a hardcoded revenue router to enforce a contractually monotone non-decreasing redemption floor. Our three contributions, stated precisely in Section~\ref{sec:introduction}, are sharper at the close. (C1) The primitive is fully specified: four immutable contracts, four free design parameters (revenue rate, LP seed fraction, asymmetric fee, treasury seed), a genesis ceremony that enforces $\Floor_0 = \Spot_0$, and a complete operation set $\mathrm{Ops}$ over which the invariants are claimed. (C2) The floor monotonicity property is proved by case analysis over $\mathrm{Ops}$ (Theorem~\ref{thm:floor-monotonicity}), the integer-arithmetic redemption residual is shown to accumulate strictly in the treasury's favor (Lemma~\ref{lem:integer-redemption}), the redemption-fairness corollary inverts the classical \citet{ChenGoldsteinJiang2010} first-mover-advantage result (Theorem~\ref{thm:redemption-fairness}), and the composite-strategy bank-run game yields a state-only saturation quantity $A^*$ above which the adversarial premium over pure redemption is a bounded constant (Theorem~\ref{thm:bank-run}). (C3) The protocol's operating envelope is characterized through a two-track simulator stack, with cross-track agreement enforced as a paper-stopping invariant and a reproducibility artifact pinning every figure to a Zenodo DOI.

The cross-domain framing we have adopted positions M\textsuperscript{2} as the on-chain implementation of three TradFi capital-return instruments that do not coexist in any single offering: the closed-end fund with NAV redemption (and our addition: monotone NAV); the perpetual sponsor-funded buyback program (with deterministic, no-announcement-gap execution and burn-not-retire disposition); and the Pigovian sell tax recycled into common-pool value rather than third-party fisc. The composition of these three instruments is the substantive design contribution; the formal-methods apparatus and the quantitative-envelope characterization are the technical contributions that support the design's defensibility.

\paragraph{What we do not claim, restated.} The work proves a structural property of the protocol's operation set. It does not eliminate smart-contract residual risk, does not promise robustness to the backing stablecoin's issuer freezing the treasury, does not address operator fraud at the revenue-deposit layer, does not predict token prices, and does not offer regulatory classification. Each of these residual risks is addressed in Section~\ref{sec:limitations} and is properly handled by deployment-time audit, counsel, and disclosure rather than by the contract code.

\paragraph{Future work.} Three threads warrant follow-up. First, an \emph{empirical} paper: after live deployment, an event-study analysis of the floor's response to revenue shocks, the redemption-rate distribution, and the realized run-resistance under genuine adversarial pressure. Such a paper would target a top finance journal (\emph{Journal of Financial Markets}, \emph{Review of Financial Studies}) and would supply the empirical content that the present mechanism-design paper cannot. Second, a \emph{mechanized-verification} paper formalizing Theorem~\ref{thm:floor-monotonicity} against the deployed bytecode under Halmos and Certora, closing open problem (P5) of Section~\ref{sec:limitations}. Third, \emph{variants}: multi-stable backing (each stable a separate token deployment; the user diversifies across instances rather than across reserves within a single instance), operator-bonding designs (the operator's revenue-deposit obligation is collateralized at genesis), and dynamic-fee schedules whose $f_s$ adjusts to the spot-floor gap. Each variant is a separate token, by the immutability discipline of the design; comparing them empirically is part of the second thread above.

The narrower technical observation we leave the reader with: the on-chain substrate makes possible a primitive---a token with a contractually-monotone-non-decreasing per-unit reserve claim---that TradFi capital-returns legal infrastructure cannot enforce. The substantive observation is that the design space for such primitives has been substantially under-explored, that the structural failure modes of prior on-chain treasury-backed tokens are not failures of the on-chain substrate but of designs that retained discretionary surface area (governance, treasury composition, post-genesis minting), and that the M\textsuperscript{2} design discipline---no discretion, no upgrade, no governance, no post-genesis minting, single transparent stable, one-way treasury inflow---is the load-bearing piece that converts a rhetorical claim into a proved theorem.
